% Apartado 2.4 - Reduccion espectral mediante FFT e IFFT (20 puntos) - BLOQUE 4
%
% Los seis incisos del enunciado:
%   a) bloques de tamano potencia de dos
%   b) preservar la simetria necesaria para reconstruir una senal real
%   c) al menos un parametro adaptativo: frecuencia interferente estimada,
%      umbral de energia, nivel de ruido o ancho de rechazo
%   d) solapamiento si la ventana o el procesamiento generan discontinuidades
%   e) aplicacion a los tres casos
%   f) al menos un caso donde eliminar mas energia del ruido degrade la senal util
%
% No se considera suficiente anular manualmente un unico bin cuya posicion se
% conozca de antemano.
% Codigo: src/filtrado_espectral.py   Figuras: fig_24_*.png
% Salidas: resultados/salidas/caso_N_espectral.wav para los tres casos

\section{Reducción espectral mediante FFT e IFFT}
\label{sec:espectral}

\textit{(pendiente — bloque 4)}

\subsection{Cadena de procesamiento por bloques}

\subsection{Estimación y modificación del espectro}

\subsection{El parámetro adaptativo}

\subsection{Reconstrucción con solapamiento}

\subsection{Aplicación a los tres casos}

\subsection{Cuándo quitar más ruido empeora el resultado}
